% Latex Document modified from Latex design of Brian R. Hall (available at www.brianrhall.net)

% Document settings
\documentclass[11pt]{article}
\usepackage[margin=1in]{geometry}
\usepackage[pdftex]{graphicx}
\usepackage{multirow}
\usepackage{setspace}
\pagestyle{plain}
\usepackage{color}
\setlength\parindent{0pt}

\begin{document}

% Course information
  \begin{center}

{\includegraphics[height=1.25in,width=1in]{wmchiffre1.jpg}} 

\LARGE Breaking Intuition: How Data has Changed Human Perception (Fall 2015, INTR 100)\\ \vspace{3mm}
\end{center}
\large \textbf{Schedule} \\
\normalsize \color{yellow}Gold \color{black} Section (CRN 14603): M, W - 12PM to 12:50PM, \textit{Morton Hall 244} \\
\color{green}Green \color{black} Section (CRN 14604): M, W - 1PM to 1:50PM, \textit{Morton Hall 244} \\
\textbf{Both} \color{black} Sections : Friday, 12:00PM to 12:50PM, \textit{Morton Hall 220} \\
\vspace{2mm}

\begin{table}[ht]
\begin{tabular}{l l}
\large \textbf{Instructors} & \\
\large Dr. \textbf{Dan} Runfola & Dr. Tyler Frazier \\
\large dsmillerrunfol@wm.edu  & tfrazier@aiddata.org \\
\large 427 Scottland Street & 427 Scottland Street \\
\large Office Hours: TBD & TBD \\
\large 757.221.1970 & 757.221.5396 \\
\end{tabular}
\end{table}

\begin{table}[ht]
\begin{tabular}{l l}
\large \textbf{Teaching Assistants} & \\
\large Dr. \textbf{Dan} Runfola & Dr. Tyler Frazier \\
\large dsmillerrunfol@wm.edu  & tfrazier@aiddata.org \\
\large 427 Scottland Street & 427 Scottland Street \\
\large Office Hours: TBD & TBD \\
\large 757.221.1970 & 757.221.5396 \\
\end{tabular}
\end{table}

\textit{Please let me know if you have any documented disabilities that may impact your performance in this class.}

\textbf {\large \\ Course Description:} The human species has nearly always sought to predict the future – from reading the clouds and sacrificial rites, to statistics and big data parsing. At the core of this history lies human intuition, and the interplay between belief, story, and knowledge. In this class, students will learn the fundamental ways in which the human species has used data over time – from arcane interpretation to artificial intelligence – and closely inspect modern day assumptions about intuition and meaning. As a part of this course, students will learn how to parse and visualize large datasets to extract meaning, and use those findings to argue for or against solutions to real-world problems. \\
\\
\textbf {Prerequisite(s):} None.

\textbf {Credit Hours:} 4 \\

\textbf {\large Books you'll need to buy:}\\ \emph{Thinking, Fast and Slow}, any edition \\
Daniel Kahneman 
\vspace{2mm}

\emph{Automate This: How Algorithms Came to Rule Our World}, any edition \\
Christopher Steiner
\vspace{2mm}

\textbf {\large Media available for free:}\\

The Beauty of Data Visualization - David McCandless \\
The Best Stats You’ve Ever Seen - Hans Rosling \\
The Rise of Human-Computer Cooperation - Shyam Sankar\\
Do schools kill creativity? - Ken Robinson \\
The Hidden Influence of Social Networks - Nicholas Christakis\\
Tableau Public - Online Software\\
R and RStudio - Free Software for Download\\
LibreOffice, Google Sheets, or Microsoft Excel\\

Throughout the course I may assign readings from academic or popular journals – these will be made available through either the Swem library or online via blackboard.

\vspace{8mm}

\textbf {\large Course Objectives:} 
\begin{enumerate} \itemsep-0.4em
  \item Challenge students traditional assumptions about the production of knowledge.
  \item Provide a more critical understanding of scientific literature.
  \item Develop your data communication, analysis, and visualization skills for future courses (and jobs).
  \item Develop critical thinking skills necessary to become more responsible for your own learning and better engage with broadly defined challenges.
  \item Provide an opportunity to fulfill the COLL 100 requirement (passing grade required).
\end{enumerate}
\vspace{8mm}

\textbf {\large Grade Distribution:} \\
\hspace*{40mm}
\begin{tabular}{ l l }
Lab Assignments & 80\% \\
Final Assignment & 10\% \\
Lecture Participation  & 10\% \\
\end{tabular} \\\\

\textbf {\large Letter Grade Distribution:} \\
\hspace*{40mm}
\begin{tabular}{ l l | l l }
\textgreater= 93.00 & A & 73.00 - 76.99 & C \\
90.00 - 92.99 & A-  & 70.00 - 72.99 & C- \\
87.00 - 89.99 & B+  & 67.00 - 69.99 & D+ \\
83.00 - 86.99 & B  & 63.00 - 66.99 & D \\
80.00 - 82.99 & B-  & 60.00 - 62.99 & D- \\
77.00 - 79.99 & C+  & \textless= 59.99 & F \\
\end{tabular} \\
.\\
% Course Policies. These are just examples, modify to your liking.
\textbf {\large Time Commitment:} Excelling in college level course work typically requires on average three to four hours per credit per week.  Since this is a four credit course, in addition to the time we meet as a class each week, you should expect to spend nine to twelve hours on average reading, writing, or otherwise preparing for this class on a weekly basis.\\

\textbf {\large Attendance:} This class does not have an attendance policy.  However, it will be difficult to learn enough to pass the class without regular participation, as the majority of course content relevant to tests and assignments will be covered in class.  Unannounced opportunities for extra credit will periodically be given in-class.\\

\textbf {\large Discussions:} Most lecture sessions will begin with a brief discussion of the assigned materials.  As such, each week students will be asked to come to class prepared with two to three bullet points for discussion - these can represent questions the material raised, commentary, or critiques.  These discussion points may be collected for extra credit.\\

\textbf {\large Classroom Behavior:} Please remain civil during discussions to promote the open exchange of ideas and foster a culture of open dialogue.  Please bear in mind that all students are entitled to their own opinion.  You are expected to listen attentively to each person speaking.  Please refrain from eating during class (and, if you must, make sure it isn't loud!).\\

\textbf {\large Teacher-student conferences:} Students performing at a C level or below are required to schedule a meeting with the instructor to discuss class performance.\\

\textbf {\large Late / Poor Performance Policy:} Assignments will not be accepted late, excepting in documented circumstances (i.e., an illness with a doctor's note).  \textit{Recognizing that it is entirely possible to underperform on key days, your lowest lab assignment grade will be dropped at the end of the semester.}\\

\textbf {\large Final Project:} The final project will be introduced the last week of classes, and will be due on the last day of finals (\textbf{December 16, 2015}). In this project, you will be assigned an open-ended question, and be asked to present a two to five page argument (including visualizations drawn from datasets) for why a certain course of action should - or should not - be taken. Example questions might include providing quantitative evidence for why you are taking the courses you have selected - or plan to select - in the next semester; contrasting the income-earning or happiness-generating potential of your planned major; or considering the likely geographic destinations for a given career path. \\

\vspace{4mm}
\textbf {\LARGE Do not cheat!} \\
.\hrulefill . \\
\textbf{Academic dishonesty is taken very seriously.  Make sure to cite all of your work, and do not turn in work that is not yours!  Cases of academic dishonesty will be evaluated and acted upon in accordance with William and Mary policies, which can be found at http://www.wm.edu/offices/deanofstudents/services/
studentconduct/} \\
.\hrulefill . \\
\vspace{8mm}

% Course Outline
\textbf {\large Course Outline}:

The weekly content might change as it depends on the progress of the class.  However, you must keep up with the reading assignments.

\begin{table}[h!]
\small % The size of the table text can be changed depending on content. Remove if desired.
\begin{tabular}{ | c | c | }
\hline
\textbf{Week} & \textbf{Content} \\
\hline
Week 1 & \begin{minipage}{.85\textwidth}
\begin{itemize} \itemsep-0.4em
	\vspace{1mm}
	\item Wed, 8/26: Intro Lab - William and Mary
	\item Fri, 8/28 - Lecture 1.0 - Preparing to be wrong  \\ \textit{Watch: Ken Robinson - Do schools kill creativity?}
	\vspace{1mm}
\end{itemize}
\end{minipage} \\
\hline
Week 2 & \begin{minipage}{.85\textwidth}
\begin{itemize} \itemsep-0.4em
	\vspace{1mm}
	\item Mon, 8/31; Wed, 9/2: Lab 1.0 - The basics I
		\item Fri, 9/4: Lecture 2.0 - What is Intuition? \\ \textit{Reading: Thinking Fast and Slow, Chapter 1 - The characters of the story}
	\vspace{1mm}
\end{itemize}
\end{minipage} \\
\hline
Week 3 & \begin{minipage}{.85\textwidth}
\begin{itemize} \itemsep-0.4em
	\vspace{1mm}
	\item Mon, 9/7; Wed, 9/9: Lab 1.1 - The basics II
	\item Fri, 9/11: Lecture 3.0 - Storytelling and Knowledge \\ \textit{Thinking Fast and Slow, Chapter 6 - Norms, Surprises and Causes} 
	 \item \textbf{Lab 1 due!}
	\vspace{1mm}
\end{itemize}
\end{minipage} \\
\hline
Week 4 & \begin{minipage}{.85\textwidth}
\begin{itemize} \itemsep-0.4em
	\vspace{1mm}
	\item Mon, 9/14; Wed, 9/16: Lab 2.0 - Deus ex machina? Considering Machine Learning I
	\item Fri, 9/18: Lecture 4.0 - The History of Modern Intuition \\ \textit{Reading: Automate This: How Algorithms Came to Rule Our World}
	\vspace{1mm}
\end{itemize}
\end{minipage} \\
\hline
Week 5 & \begin{minipage}{.85\textwidth}
\begin{itemize} \itemsep-0.4em
	\vspace{1mm}
	\item Mon, 9/21; Wed, 9/23: Lab 2.1 - Deus ex machina? Considering Machine Learning II
	\item \textbf{Fri, 9/25: Lecture 5.0: Guest Lecturer TBD} 
	\item \textbf{Lab 2 due!}
	\vspace{1mm}
\end{itemize}
\end{minipage} \\
\hline
Week 6 & \begin{minipage}{.85\textwidth}
\begin{itemize} \itemsep-0.4em
	\vspace{1mm}
	\item Mon, 9/28; Wed, 9/30: Lab 3.0 - Data Processing and Social Data I
	\item Fri, 10/2: Lecture 6.0 - Applications of Knowledge \\ \textit{The Hidden Influence of Social Networks - Nicholas Christakis}
	\vspace{1mm}
\end{itemize}
\end{minipage} \\
\hline
Week 7 & \begin{minipage}{.85\textwidth}
\begin{itemize} \itemsep-0.4em
	\vspace{1mm}
	\item Mon, 10/5; Wed, 10/7: Lab 3.1 - Data Processing and Social Data II
	\item \textbf{Fri, 10/9: Lecture 7.0: Guest Lecturer TBD}
	\vspace{1mm}
\end{itemize}
\end{minipage} \\
\hline
Week 8 & \begin{minipage}{.85\textwidth}
\begin{itemize} \itemsep-0.4em
	\vspace{1mm}
	\item Mon, 10/12: No Class (Fall Break)
	\item Wed, 10/14: Lab 3.2 - Creating Powerful Visualizations
	\item Fri, 10/16: Lecture 8.0: Data Visualization: Perception and Knowledge  \\ \textit{The Beauty of Data Visualization - David McCandless} 
	\item \textbf{Lab 3 due!}
	\vspace{1mm}
\end{itemize}
\end{minipage} \\
\hline
Week 9 & \begin{minipage}{.85\textwidth}
\begin{itemize} \itemsep-0.4em
	\vspace{1mm}
	\item \textbf{Mon, 10/19; Wed 10/21: Lab 4.0 - Getting Rich Quick: Intuition vs. Machines on the Stock Market (*Placeholder for Business Analytics)}
	\item Fri, 10/23: Lecture 9.0: The meaning of Big Data  \\ \textit{Automate This: How Algorithms Came to Rule Our World, Chapter 1: Wall Street, the First Domino}
	\vspace{1mm}
\end{itemize}
\end{minipage} \\
\hline
Week 10 & \begin{minipage}{.85\textwidth}
\begin{itemize} \itemsep-0.4em
	\vspace{1mm}
		\item \textbf{Mon, 10/26; Wed 10/28: Lab 4.1 - Getting Rich Quick: Intuition vs. Machines on the Stock Market (*Placeholder for Business Analytics)} 
	\item \textbf{Fri, 10/30: Lecture 10.0: Guest Lecturer TBD, likely D. Murray}
	\item \textbf{Lab 4 due!}
	\vspace{1mm}
\end{itemize}
\end{minipage} \\
\hline
Week 11 & \begin{minipage}{.85\textwidth}
\begin{itemize} \itemsep-0.4em
	\item Mon, 11/2; Wed 11/4: Lab 5.0 - Changing Behavior from Outer Space I
	\item Fri, 11/6: Lecture 11.0: Why we're generally wrong  \\ \textit{Thinking, Fast and Slow - Chapter 18: Taming Intuitive Predictions}
	\vspace{1mm}
\end{itemize}
\end{minipage} \\
\hline
Week 12 & \begin{minipage}{.85\textwidth}
\begin{itemize} \itemsep-0.4em
	\vspace{1mm}
	\item Mon, 11/9; Wed 11/11: Lab 5.1 - Changing Behavior from Outer Space II
	\item \textbf{Fri, 11/13: Lecture 12.0: Guest Lecturer TBD, likely R. Rose}
	\item \textbf{Lab 5 due!}
	\vspace{1mm}
\end{itemize}
\end{minipage} \\
\hline
Week 13 & \begin{minipage}{.85\textwidth}
\begin{itemize} \itemsep-0.4em
	\vspace{1mm}
	\item Mon, 11/16; Wed 11/18: Lab 6.0 - Digging into the Past Using Data I
	\item \textbf{Fri, 11/20: Lecture 13.0: Final Guest Lecture}
	\vspace{1mm}
\end{itemize}
\end{minipage} \\
\hline
Week 14 & \begin{minipage}{.85\textwidth}
\begin{itemize} \itemsep-0.4em
	\vspace{1mm}
	\item Mon, 11/23: Lab 6.1 - Digging into the Past Using Data II
	\item Wed, 11/15; Fri, 11/27 - No Class (Thanksgiving Break)
	\vspace{1mm}
\end{itemize}
\end{minipage} \\
\hline
Week 15 & \begin{minipage}{.85\textwidth}
\begin{itemize} \itemsep-0.4em
	\vspace{1mm}
	\item Mon, 11/30, Wed 12/2: Final Project Introduction and Work Time \textbf{(Due 12/16)}
	\item Fri, 12/4: Lecture 14.0: Breaking Intuition  \\ \textit{Thinking, Fast and Slow - Chapter 36: Life as a Story;\\ Automate This - Chapter 10 - The Future} 
	\item \textbf{Lab 6 due!}
	\vspace{1mm}
\end{itemize}
\end{minipage} \\
\hline
\end{tabular} 
\end{table}

\end{document}



